\chapter{Introduction}\label{introduction}
% \textcolor{red}{The introduction of your bachelor thesis introduces the research area, the research hypothesis, and the scientific contributions of your work.
% A good narrative structure is the one suggested by Simon Peyton Jones~\cite{peys04:HowToWriteAGoodResearchPaper}:}

% \begin{itemize}
% 	\item \textcolor{red}{describe the problem \slash{} research question}
% 	\item \textcolor{red}{motivate why this problem must be solved}
% 	\item \textcolor{red}{demonstrate that a (new) solution is needed}
% 	\item \textcolor{red}{explain the intuition behind your solution}
% 	\item \textcolor{red}{motivate why \slash{} how your solution solves the problem (this is technical)}
% 	\item \textcolor{red}{explain how it compares with related work}
% \end{itemize}

% \textcolor{red}{Close the introduction with a paragraph in which the content of the next chapters is briefly mentioned (one sentence per chapter).}



% - TALK ABOUT RISE OF LLM USAGE SINCE 2020 - introduction of problem
In recent years, large language models (LLMs) have emerged and evolved rapidly \cite{zhao2023survey}. While initially restricted to research environments, these tools have now reached massive public scale. For instance, OpenAI recently confirmed it has surpassed 50 million paying subscribers across its commercial tiers, with up to 900 million weekly active users processing over 2 billion prompts per day \cite{openai2025enterprise}. This scale demonstrates that interacting with LLMs has become a routine global practice. However, while these models assist users in learning, working, and automating tasks, their widespread deployment does not come without downsides.

One primary risk is \textit{hallucination}, a phenomenon further discussed in Chapter \ref{preliminaries}, where an LLM generates output that is unfaithful to its original source or the user's instructions. While minor discrepancies often go unnoticed, severe hallucinations can have profound consequences. For example, Zhao \etal \cite{zhoa2026llmhallucinationsinthewild} analyzed 2.5 million published papers and found a sharp rise in non-existent references following widespread LLM adoption, conservatively estimating over 146,000 hallucinated citations in 2025 alone. Similar impacts have been documented across other critical sectors. In healthcare, hallucinations pose severe ethical and safety risks \cite{healthcarehallucinations2025}, while in the commercial sector, the use of ungrounded conversational agents has led to concrete operational liabilities for enterprises \cite{lseo2026liability}.

As mentioned above, the occurrences of hallucinations are well-documented in high-stakes fields like medicine and law. 
In this thesis, however, we will consider a use-case which is less universally known, a case on a local music venue in Nijmegen. 
Since local venues often give the stage to smaller, lesser-known artists, the likelihood of hallucinations increases. While large language models demonstrate impressive factual recall for globally prominent topics, their reliability degrades significantly when processing niche or localized information. This limitation, often referred to as popularity bias, is formally quantified by Kandpal \etal \cite{kandpal2023longtail}, who demonstrate that an LLM's ability to accurately state a fact is directly correlated to the frequency of that entity's appearance in its pre-training corpus.

To address these widespread hallucinations, a common modern approach is the ``LLM-as-a-judge'' solution, wherein a highly capable model (such as GPT-4) is prompted to verify the factual accuracy of generated text. However, recent benchmarking reveals critical structural flaws in this methodology. Chen \etal \cite{chen2023felm} introduced the FELM benchmark, demonstrating that LLM evaluators struggle significantly to detect factual errors, often failing to identify hallucinations or confidently hallucinating false corrections themselves. This failure occurs because an LLM acting as a judge relies on the same flawed, popularity-biased parametric memory as the model generating the text, where a parametric memory refers to a models internal knowledge determined by the models internal knowledge. If an LLM lacks the training data to accurately describe a local Dutch punk band, it also lacks the data required to fact-check a summary about that band. Consequently, relying on an LLM to evaluate long-tail cultural metadata creates a circular validation flaw.

To solve this, a new approach to automated evaluation is required. Traditional hallucination benchmarks often evaluate models against static, global datasets and primarily focus purely on English tasks. 
These traditional measures where not set up with testing small scale musical knowledge in mind, where the source text differs from the summarisation. 

Therefore, we propose a two-phase hybrid evaluation architecture designed specifically for this domain.

The intuition behind this solution is to split the evaluation into two distinct steps: confirming the source text was understood, and checking if the claims match reality. First, to assess whether the LLM faithfully translated and summarized the provided text, we utilize a cross-lingual Natural Language Inference (NLI) model (XLM-RoBERTa). This isolates intrinsic hallucinations. Second, to overcome the LLM's popularity bias regarding long-tail entities, we integrate a non-parametric, structured external database (MusicBrainz). By deterministically cross-referencing the LLM's claims against this real-world knowledge graph, we can catch extrinsic hallucinations that the NLI model alone would miss. Using manually labelled claims, we compare the accuracy of our built pipeline and are able to identify wide range of hallucinations types. 


Our primary research question entails: 
\begin{enumerate}
    \item[\textbf{RQ 1}] To what extent can a hybrid evaluation framework, integrating deterministic claim decomposition, cross-lingual Natural Language Inference (NLI), and a structured external knowledge graph, effectively detect and categorize hallucinations in one-sentence LLM-generated local music venue summaries?
    \begin{enumerate}
        \item[\textbf{RQ 1.1}] How accurately can the cross-lingual NLI baseline pipeline detect and classify intrinsic hallucinations by evaluating semantic entailment between English atomic claims and their Dutch source texts, and to what extent does it label extrinsic hallucinations with a neutral or supported verdict?
        \item[\textbf{RQ 1.2}] How does the calibration of the baseline confidence threshold ($\mathcal{T}$) impact the models error detection and accuracy? 
        \item[\textbf{RQ 1.3}] How effectively can the integration of the MusicBrainz external knowledge graph increase the accuracy of atomic and sentence level labels of the baseline by validating extrinsic and intrinsic hallucinations?
        \item[\textbf{RQ 1.4}] How does aggregating these atomic claim-level evaluations using min-pooling impact the  error detection and false positive rates at the sentence level?
    \end{enumerate}
\end{enumerate}

The remainder of this thesis is structured as follows. Chapter \ref{preliminaries} outlines the theoretical background regarding hallucinations in LLMs, cross-lingual NLI, Knowledge Graphs and Binary Classification Metrics. Next, Chapter \ref{Newsletter} discusses the used dataset, how it was generated and what existing problems occur. Following this, Chapter \ref{research} discusses the full evaluation pipeline architecture, including generating the corpus, extracting source data and the external knowledge graph, parsing Dutch source text, decomposing the summaries into atomic claims, validating and evaluating these atomic claims and lastly the experiment set up. Chapter \ref{results} shows the results of our experiments, which will then be discussed in Chapter \ref{discussion} and future work will be suggested in addition to our limitations . Finally, Chapter \ref{conclusions} concludes the research by answering the research questions according to our results and discussion.




